\section{Manipulación robótica}

\subsection{Los retos de la manipulación diestra}

La manipulación diestra (por ejemplo, sujetar y manipular objetos con una mano de varios dedos) es más difícil que la locomoción:
\begin{itemize}
    \item La dinámica de contacto es fuertemente no lineal
    \item El espacio de estados tiene más dimensiones
    \item Las tareas son más diversas
    \item La Sim-to-real gap es más grave en el modelado del contacto
\end{itemize}

\subsection{Aplicaciones del RL en la manipulación}

El ponente presentó varios casos de manipulación diestra con RL, entre ellos:
\begin{itemize}
    \item Girar un cubo de Rubik (OpenAI Rubik's Cube)
    \item Hacer girar un bolígrafo (pen spinning)
    \item Sujeción diestra (dexterous grasping)
\end{itemize}

Todas estas tareas siguen el mismo enfoque: simulación paralela a gran escala, domain randomization y teacher-student.

\subsection{Resumen de la sección}

El RL ya ha mostrado una capacidad superior a la humana en diversas tareas de manipulación robótica, pero la calidad del simulador sigue siendo un factor limitante decisivo.

